\chapter*{Introducere}
\addcontentsline{toc}{chapter}{Introducere}

Sectorul financiar-bancar traversează în prezent cea mai amplă și accelerată transformare tehnologică din istoria sa modernă. Migrarea accelerată de la modelele bancare tradiționale, bazate pe sucursale fizice și canale de comunicare închise, către ecosisteme financiare digitale descentralizate, platforme de open banking, microservicii în cloud și interfețe de programare a aplicațiilor (API-uri) a redefinit fundamental modul în care instituțiile de credit interacționează cu clienții și cu partenerii de afaceri. Această evoluție a adus beneficii economice substanțiale: tranzacții în timp real, costuri operaționale optimizate, accesibilitate continuă și o eficiență sporită a fluxurilor de capital.

Totuși, digitalizarea extinsă a condus la o expansiune fără precedent a suprafeței de atac. Pentru entitățile rău-intenționate – variind de la grupări de criminalitate cibernetică organizată motivate financiar, până la actori statali persistenți avansați (APT – Advanced Persistent Threats) – sistemele informatice bancare reprezintă ținta supremă. O breșă de securitate într-o instituție financiară nu produce doar pierderi patrimoniale directe, ci poate declanșa crize de lichiditate, falimente bancare, erodarea ireversibilă a încrederii publice și chiar instabilitate sistemică la nivel macroeconomic. În acest peisaj operațional ostil, abordările convenționale de securitate, bazate exclusiv pe apărarea perimetrică clasică, s-au dovedit complet inadecvate.

Conștientizând această vulnerabilitate sistemică, organismele internaționale și europene de reglementare au instituit cadre normative extrem de riguroase. Directiva Revizuită privind Serviciile de Plată (PSD2) a impus mecanisme obligatorii de Autentificare Strictă a Clienților (SCA – Strong Customer Authentication) și legături dinamice tranzacționale. Standardul de Securitate a Datelor din Industria Cardurilor de Plată (PCI-DSS v4.0) a stabilit cerințe stricte pentru criptarea și izolarea mediului datelor titularilor de card (CDE). În paralel, noul Regulament European privind Reziliența Operațională Digitală (DORA – Regulamentul UE 2022/2554) obligă instituțiile financiare să demonstreze nu doar existența unor controale defensive, ci capacitatea demonstrabilă de a rezista, de a absorbi șocurile și de a se recupera rapid în urma unor incidente cibernetice disruptive majore.

Alegerea acestei teme de cercetare este motivată de necesitatea acută de a crea o punte solidă între teoria economică financiar-bancară (evidența contabilă în partidă dublă, decontarea tranzacțiilor, gestiunea riscurilor) și ingineria sistemelor informatice securizate (rețele de calculatoare, sisteme de operare întărite, baze de date tranzacționale și securitate ofensivă/defensivă). În cadrul programului de studii Informatică Economică din cadrul Facultății de Economie și Administrarea Afacerilor a Universității din Craiova, înțelegerea mecanismelor interne care garantează integritatea și confidențialitatea activelor digitale reprezintă o competență fundamentală pentru formarea specialiștilor capabili să gestioneze riscul operațional modern.

Obiectivul principal al prezentei lucrări de licență constă în proiectarea, implementarea completă, testarea experimentală și auditarea unei arhitecturi informatice bancare reziliente, complet virtualizate într-un laborator de cercetare avansat. Pentru îndeplinirea acestui deziderat general, au fost stabilite următoarele obiective specifice de cercetare:

• Obiectivul 1: Analiza critică a stadiului actual al cunoașterii privind securitatea cibernetică în mediul bancar, investigând cerințele de conformitate (DORA, PSD2 RTS, PCI-DSS v4.0, ISO 20022), principiile contabile ale sistemelor Core-Banking și vectorii de atac contemporani.

• Obiectivul 2: Proiectarea unei topologii de laborator complet virtualizate pe platforma de hypervisor bare-metal Proxmox VE 9.2, bazată pe segmentare logică de rețea prin VLAN-uri (Management, Services, CyberLab) și firewalling perimetral de ultimă generație prin OPNsense.

• Obiectivul 3: Implementarea practică a celor patru active bancare fundamentale: un motor central de Core-Banking (VM 310) inspirat de standardul deschis Apache Fineract cu registru contabil în partidă dublă și înlănțuire criptografică SHA-256 a tranzacțiilor; un server de bază de date relațională (VM 311 - PostgreSQL) cu tabele de audit nealterabile; o poartă de plăți (CT 312) cu validare Luhn, verificare anti-fraudă și suport pentru schemele de mesagerie SWIFT MT103 și ISO 20022 pacs.008; precum și un bastion de acces securizat (VM 313 - Jump-Box) protejat prin chei criptografice Ed25519 și autentificare cu factori multipli (MFA).

• Obiectivul 4: Dezvoltarea și executarea unei suite automatizate de simulare a atacurilor cibernetice (Red Team vs. Blue Team) acoperind 5 scenarii critice: de la tranzacționarea nominală de referință, la atacuri de tip SQL Injection și alterare directă de sold (Balance Tampering), atacuri de tip Card Stuffing pe fluxul de autorizare, mișcare laterală și eludare a bastionului administrativ, până la coruperea structurilor de mesagerie interbancară.

• Obiectivul 5: Integrarea unui sistem de monitorizare a integrității și detecție a intruziunilor (HIDS/SIEM - Wazuh), configurarea unor reguli analitice de corelare în timp real capabile să identifice decalajele contabile și maparea integrală a incidentelor pe matricea internațională MITRE ATT\&CK for Financial Services.

• Obiectivul 6: Evaluarea riguroasă a performanței operaționale, a rezilienței la incidente și a gradului de conformitate de reglementare obținut, evidențiind valoarea adăugată a cercetării și impactul economic direct asupra diminuării pierderilor operaționale din fraude.

Metodologia de cercetare adoptată îmbină analiza teoretică și normativă a literaturii de specialitate cu o abordare cantitativă experimentală. Întreaga arhitectură a fost instanțiată pe nodul fizic de calcul x86\_64 bare-metal al infrastructurii de laborator (procesor Intel Core i3-10100F cu 4 nuclee fizice / 8 fire de execuție la 4.30 GHz Turbo, 12 GB memorie RAM DDR4 completată cu subsistem dinamic ZRAM de 6.0 GB lz4 și subsistem de stocare de 512 GB SSD LVM-Thin sub Proxmox VE 9.2). Toate componentele software au fost dezvoltate în limbajul Python 3, utilizând cadre de lucru moderne (FastAPI, SQLite/PostgreSQL, hashlib) și au fost supuse unor teste de penetrare automate riguroase, asigurând reproductibilitatea completă a rezultatelor obținute.

Lucrarea este structurată riguros pe două capitole principale, conform ghidului de elaborare al facultății. Primul capitol, „Stadiul cunoașterii în securitatea cibernetică bancară”, sintetizează fundamentele legislative, arhitecturale, vulnerabilitățile curente și paradigmele defensive de ultimă generație. Al doilea capitol, „Proiectarea, implementarea și evaluarea arhitecturii bancare reziliente într-un mediu virtualizat”, constituie contribuția originală extinsă a autorului, detaliind construcția mediului experimental, codul sursă al serviciilor bancare, derularea atacurilor controlate, telemetria sistemului SIEM și validarea ipotezelor de cercetare. Lucrarea se încheie cu o secțiune de concluzii, bibliografia consultată și anexe tehnice cuprinzătoare.

